\section{Kết quả}
\label{sec:results}

Toàn bộ bảng và trị số thống kê trong phần này được sinh tự động từ artifact kết quả
bằng pipeline phân tích (\cref{sec:reproducibility}), nhằm bảo đảm số trong bảng và
trong prose luôn khớp với artifact. \TBD{Các lượt chạy trên dải Qwen đang được thực
hiện; mọi số liệu hiển thị hiện là placeholder ở chế độ demo và sẽ được thay bằng số
thực khi pipeline hoàn tất.}

\subsection{RQ1: Chất lượng theo model scale}
\label{sec:rq1}

\Cref{tab:rq1-ladder} trình bày Total của bảy tác tử trên \texttt{test\_90} theo bốn
kích thước Qwen; mỗi ô là trung bình per-scenario gộp ba lượt chạy độc lập, kèm độ
lệch chuẩn giữa các lượt và $n$ scenario complete-case. Ở kích thước lớn nhất (9B),
AlignmentGraph-ISD đạt Total $\rqOneHarnessTotalNineB$
(SD giữa các lượt $\rqOneHarnessRunSDNineB$), so với $\rqOneBestBaselineTotalNineB$
của đối sánh mạnh nhất tại kích thước này (\rqOneBestBaselineNameNineB); chênh lệch
ghép cặp tương ứng là $+\rqOneDeltaBestNineB$ điểm với khoảng tin cậy bootstrap 95\%
$[\rqOneDeltaBestCILoNineB;\,\rqOneDeltaBestCIHiNineB]$. Ở kích thước nhỏ nhất
(0,8B), harness đạt $\rqOneHarnessTotalZeroEightB$ so với
$\rqOneBestBaselineTotalZeroEightB$ của \rqOneBestBaselineNameZeroEightB{}
(chênh lệch $+\rqOneDeltaBestZeroEightB$;
CI $[\rqOneDeltaBestCILoZeroEightB;\,\rqOneDeltaBestCIHiZeroEightB]$). Trên toàn bộ
24 so sánh ghép cặp (bốn kích thước, sáu đối sánh mỗi kích thước, hiệu chỉnh Holm
trong từng kích thước), $p_{\mathrm{Holm}}$ lớn nhất là $\rqOnePHolmMaxAll$.
\TBD{prose bổ sung: thứ hạng đầy đủ theo từng kích thước, các đối sánh suy giảm mạnh,
và nhận xét về SD giữa các lượt, khi có số thực}.

% AUTO-GENERATED by isd-agent-benchmark/scripts/alignmentgraph-isd-bench/09_gen_paper_tables.py — DO NOT EDIT BY HAND
% tab_rq1_ladder.tex — RQ1 ladder table on ADDIE
% generated_at: 2026-07-30T17:41:41
% source pooled_ladder.json generated_at: 2026-07-30T15:54:18

\begin{table*}[t]
  \caption{RQ1: ISD-Agent-Bench \texttt{test\_90} theo thang kích thước mô hình (Qwen, 3 run độc lập mỗi kích thước). Mỗi ô: trung bình rubric ADDIE per-scenario (/100) — thước đo chất lượng nội dung thiết kế giảng dạy, tức construct mà câu hỏi nghiên cứu nhắm tới, gộp mean-of-runs trên các scenario complete-case; giá trị sau $\pm$ là SD giữa 3 run (run-to-run); $n$ = số scenario đủ cả 3 run. \textbf{Đậm} = tốt nhất theo cột.}
  \label{tab:rq1-ladder}
  \small
  \begin{tabular}{@{}lr@{\hskip 3pt}rr@{\hskip 3pt}rr@{\hskip 3pt}rr@{\hskip 3pt}r@{}}
    \toprule
    & \multicolumn{2}{c}{Qwen3.5-0,8B} & \multicolumn{2}{c}{Qwen3.5-2B} & \multicolumn{2}{c}{Qwen3.5-4B} & \multicolumn{2}{c}{Qwen3.5-9B} \\
    \cmidrule(lr){2-3}\cmidrule(lr){4-5}\cmidrule(lr){6-7}\cmidrule(lr){8-9}
    Agent & ADDIE {\scriptsize$\pm$SD} & $n$ & ADDIE {\scriptsize$\pm$SD} & $n$ & ADDIE {\scriptsize$\pm$SD} & $n$ & ADDIE {\scriptsize$\pm$SD} & $n$ \\
    \midrule
    Baseline & 44{,}52{\scriptsize$\pm$1{,}19} & 5\textsuperscript{*} & 57{,}82{\scriptsize$\pm$2{,}01} & 47\textsuperscript{*} & 78{,}68{\scriptsize$\pm$0{,}48} & 82\textsuperscript{*} & 83{,}45{\scriptsize$\pm$0{,}32} & 86\textsuperscript{*} \\
    EduPlanner & 46{,}61{\scriptsize$\pm$2{,}60} & 7\textsuperscript{*} & 55{,}12{\scriptsize$\pm$0{,}58} & 57\textsuperscript{*} & 77{,}61{\scriptsize$\pm$0{,}27} & 77\textsuperscript{*} & 80{,}31{\scriptsize$\pm$0{,}70} & 85\textsuperscript{*} \\
    ADDIE-Agent & 73{,}55{\scriptsize$\pm$0{,}13} & 90 & 76{,}98{\scriptsize$\pm$0{,}20} & 90 & 79{,}82{\scriptsize$\pm$0{,}05} & 90 & 80{,}42{\scriptsize$\pm$0{,}06} & 90 \\
    RPISD-Agent & 77{,}79{\scriptsize$\pm$0{,}11} & 88\textsuperscript{*} & 79{,}69{\scriptsize$\pm$0{,}20} & 90 & 81{,}03{\scriptsize$\pm$0{,}17} & 90 & 81{,}48{\scriptsize$\pm$0{,}06} & 90 \\
    Dick-Carey-Agent & 79{,}74{\scriptsize$\pm$0{,}31} & 90 & 79{,}62{\scriptsize$\pm$0{,}26} & 90 & 81{,}84{\scriptsize$\pm$0{,}34} & 89\textsuperscript{*} & 78{,}88{\scriptsize$\pm$0{,}13} & 90 \\
    ReAct-ADDIE & 77{,}48{\scriptsize$\pm$0{,}36} & 90 & 82{,}19{\scriptsize$\pm$0{,}25} & 90 & 85{,}08{\scriptsize$\pm$0{,}15} & 90 & \textbf{86{,}30}{\scriptsize$\pm$0{,}08} & 90 \\
    \textbf{AlignmentGraph-ISD} & \textbf{80{,}46}{\scriptsize$\pm$0{,}26} & 90 & \textbf{84{,}69}{\scriptsize$\pm$0{,}02} & 90 & \textbf{85{,}88}{\scriptsize$\pm$0{,}16} & 90 & 86{,}23{\scriptsize$\pm$0{,}15} & 90 \\
    \bottomrule
  \end{tabular}
  \par\smallskip\footnotesize\textsuperscript{*}$n<$ tổng số scenario: có scenario lỗi/thiếu ở ít nhất một run (complete-case; xem stats\_summary.md).
\end{table*}


\paragraph{Interaction theo kích thước mô hình.} Để kiểm tra nhận định về khoảng cách
theo kích thước bằng một phép kiểm định trực tiếp thay vì so hai ước lượng riêng,
chúng tôi báo difference-in-differences của chênh lệch harness trừ Baseline giữa kích
thước lớn nhất và nhỏ nhất: $\rqOneDiDVsBaseline$ điểm (CI 95\%
$[\rqOneDiDVsBaselineCILo;\,\rqOneDiDVsBaselineCIHi]$;
$p=\rqOneDiDVsBaselineP$), cùng độ dốc xu hướng của chênh lệch theo bậc kích thước
($\rqOneTrendSlopeVsBaseline$ điểm/bậc). Kết luận về việc lợi thế của harness có nới
rộng khi mô hình nhỏ đi hay không dựa trên dấu và khoảng tin cậy của các đại lượng
này. \TBD{diễn giải interaction khi có số thực}.

\subsection{RQ2: Constructive alignment theo tiêu chí đã thiết lập}
\label{sec:rq2}

\Cref{tab:rq2-alignment} trình bày các endpoint cấp objective
(\cref{sec:alignment-metric})---full alignment rate (FAR, endpoint chính), assessment
coverage, activity coverage, cognitive congruence---cùng alignment composite cho bảy
tác tử theo bốn kích thước mô hình. Chúng tôi báo cáo ba lớp kết quả: (i) so sánh
ghép cặp giữa harness và từng đối sánh trong từng kích thước trên FAR và composite
(Wilcoxon signed-rank, hiệu chỉnh Holm với family sáu so sánh cho từng chỉ số,
khoảng tin cậy bootstrap); (ii) sự thay đổi của từng endpoint và thành phần khi kích
thước mô hình giảm; và (iii) interaction method~$\times$~model size trên composite
bằng difference-in-differences, kiểm tra trực tiếp giả thuyết rằng khoảng cách
alignment nới rộng khi mô hình nhỏ đi. \TBD{prose kết quả RQ2: mức FAR và composite,
chênh lệch ghép cặp, CI, $p_{\mathrm{Holm}}$ và kết quả interaction, khi có số thực}.

% AUTO-GENERATED by isd-agent-benchmark/scripts/alignmentgraph-isd-bench/09_gen_paper_tables.py — DO NOT EDIT BY HAND
% tab_rq2_alignment.tex — RQ2 objective->assessment similarity x model sizes
% generated_at: 2026-07-30T17:41:41
% source pooled_ladder.json generated_at: 2026-07-30T15:54:18

\begin{table}[t]
  \caption{RQ2: objective$\to$assessment similarity (mean-max rectified cosine, liên tục, không ngưỡng) --- vế \emph{measures} của \textit{aligned(o)}; trung bình per-scenario (SD), gộp mean-of-runs, theo kích thước mô hình. Sáu tín hiệu còn lại của panel: Bảng~\ref{tab:rq2-panel}.}
  \label{tab:rq2-alignment}
  \small
  \begin{tabular}{@{}lrrrr@{}}
    \toprule
    Agent & Qwen3.5-0,8B & Qwen3.5-2B & Qwen3.5-4B & Qwen3.5-9B \\
    \midrule
    Baseline & 0{,}104 (0{,}20) & 0{,}383 (0{,}17) & 0{,}547 (0{,}09) & 0{,}562 (0{,}05) \\
    EduPlanner & 0{,}360 (0{,}23) & 0{,}500 (0{,}16) & 0{,}598 (0{,}09) & 0{,}592 (0{,}07) \\
    ADDIE-Agent & 0{,}595 (0{,}06) & 0{,}606 (0{,}08) & 0{,}608 (0{,}06) & 0{,}586 (0{,}06) \\
    RPISD-Agent & 0{,}510 (0{,}07) & 0{,}519 (0{,}05) & 0{,}466 (0{,}04) & 0{,}451 (0{,}04) \\
    Dick-Carey-Agent & 0{,}441 (0{,}05) & 0{,}499 (0{,}05) & 0{,}536 (0{,}05) & 0{,}557 (0{,}05) \\
    ReAct-ADDIE & 0{,}504 (0{,}11) & 0{,}583 (0{,}10) & 0{,}579 (0{,}06) & 0{,}570 (0{,}06) \\
    \textbf{AlignmentGraph-ISD} & 0{,}851 (0{,}05) & 0{,}850 (0{,}06) & 0{,}806 (0{,}07) & 0{,}832 (0{,}05) \\
    \bottomrule
  \end{tabular}
\end{table}


Vì quy trình đánh giá đọc cùng một văn bản kết xuất cho mọi tác tử và không dùng phán
đoán LLM ở bất kỳ bước tổng hợp nào, kết quả RQ2 là một phép đo tự động của vị từ
$\mathrm{aligned}(o)$ trong \cref{eq:aligned} theo các tiêu chí đã thiết lập, độc lập
với cả verifier của harness lẫn LLM judge của benchmark; các so sánh được diễn giải
ở mức tương đối giữa các tác tử. Mối liên hệ giữa các chỉ số alignment và điểm judge
của RQ1---hai phép đo nhắm vào hai khía cạnh khác nhau của chất lượng---được thảo
luận trong \cref{sec:discussion}.
